% Part: computability
% Chapter: computability-theory
% Section: no-universal-function

\documentclass[../../../include/open-logic-section]{subfiles}

\begin{document}

\olfileid{cmp}{thy}{nou}
\olsection{لا توجد دالة شاملة قابلة للحساب}

مع أن ثمة دالة جزئية قابلة للحساب شاملة للدوال الجزئية القابلة للحساب،
فلا توجد دالة كلية قابلة للحساب شاملة للدوال الكلية القابلة للحساب.

\begin{thm}
\ollabel{thm:no-univ}
لا توجد دالة شاملة قابلة للحساب. وبعبارة أخرى، إذا كانت الدالة
$\fn{Un}'(k,x)$ تحقق أنه لكل دالة كلية قابلة للحساب $f(x)$ عدد طبيعي
~$k$ بحيث $f(x)=\fn{Un}'(k,x)$ لكل~$x$، فإن تلك الدالة ليست قابلة
للحساب.
\end{thm}

\begin{proof}
البرهان حجة قطرية بسيطة: لو كانت $\fn{Un}'(k,x)$ كلية وقابلة للحساب،
لكانت
\[
d(x) = \fn{Un}'(x, x) + 1
\]
كلية وقابلة للحساب أيضًا. غير أن $d(k)$، بحسب التعريف، لا تساوي
$\fn{Un}'(k,k)$. ومن ثم تختلف قيمتا $d(x)$ و$\fn{Un}'(k,x)$، لكل
$k$، عند~$x$ واحد على الأقل، وهو $x=k$.
\end{proof}

\begin{explain}
تبين \olref[uni]{thm:univ-comp} أعلاه أننا نستطيع الإفلات من هذه الحجة
القطرية، ولكن بثمن السماح للدالة الشاملة بأن تكون جزئية. أي إن
$\fn{Un}$ شاملة للدوال الكلية القابلة للحساب، لكنها ليست كلية. ولا
تنجح الحجة القطرية في الحالة الجزئية.
\end{explain}

\begin{prob}
  لفهم سبب عدم نجاح الحجة القطرية في برهان \olref{thm:no-univ} في
  الحالة الجزئية، انظر في الدالة
  $f(x)\simeq\fn{Un}(x,x)+1$. هل هي جزئية قابلة للحساب؟ إذا كانت كذلك،
  فلها فهرس~$e$، أي $f(x)\simeq\fn{Un}(e,x)$. فماذا يمكنك أن تقول عن
  ~ $f(e)$؟
\end{prob}


\end{document}
